\GuideSpread{What the coordinates must preserve}{From an integer source to a strict descent certificate}{01_route.pdf}{
The top object is a positive odd integer. Its valuations come from the actual orbit. They are not independent choices.\par
The ordered word determines a dyadic depth $S$ and an offset $C$. Together they retain the exact endpoint.\par
The bottom inequality is the test for falling below the original source. Each diagram in this guide develops one part of that chain.\par
An exact certificate can apply to one seed or to an infinite arithmetic progression. A global proof must force an applicable certificate for every odd seed above 1.
}{A useful coordinate system preserves the original source and the entire offset.}

\GuideSpread{Three maps, three clocks}{One orbit from 3, counted three different ways}{02_clocks.pdf}{
The classical map applies $3n+1$ at an odd input and divides by two at an even input.\par
The shortcut map performs the first division by two together with an odd multiplication.\par
The odd map $A$ removes all factors of two after that multiplication. Its states are always odd.\par
All three rows reach 1. They use 7, 5, and 2 steps respectively. This distinction matters whenever a computation or a formal certificate reports an iteration count.
}{A step count is meaningful only after the map has been named.}

\GuideSpread{Valuation is an exact integer measurement}{How many factors of two occur in $3n+1$?}{03_valuations.pdf}{
For an odd source, $a(n)=\nu_2(3n+1)$. It is the largest exponent of two dividing the numerator.\par
For example, $3\cdot5+1=16$ gives $a(5)=4$ and $A(5)=1$. The source 7 instead has $3\cdot7+1=22$, so $a(7)=1$ and $A(7)=11$.\par
The bars show every positive odd source at most 31. They are exact integer values. A high valuation may cause a strong fall, but a finite collection of such examples cannot settle every orbit.
}{The division exponent belongs to the actual state.}

\GuideSpread{The affine identity is a balance}{The two step word $(1,2)$ from 27}{04_affine.pdf}{
The states are $27,41,31$. Two steps accumulate depth $S=3$. The ordered offset is $C=5$.\par
The equality is $8\cdot31=9\cdot27+5$. The first term on the right is the homogeneous transfer. The gold end of its bar is the offset.\par
Every realized word obeys $2^Sx_k=3^kn+C$. This is an exact integer identity. Its two parts must remain together when the endpoint is compared with the source.
}{The small offset can decide the sign of a marginal comparison.}

\GuideSpread{Order survives in the offset}{Equal length and depth do not give equal transfers}{05_order.pdf}{
The words $(1,3)$ and $(3,1)$ both have length 2 and depth 4. Their offsets are 5 and 11.\par
To compute them, start at $S_0=C_0=0$ and use $C_{i+1}=3C_i+2^{S_i}$. Earlier divisions change the weight of the later additive terms.\par
The two words also have different exact integer cylinders. The formula for a transfer is evaluated on its own cylinder. An average valuation would identify these words and discard both differences.
}{An ordered packet contains more information than its total depth.}

\GuideSpread{The endpoint needs one more bit}{Integrality alone does not certify the word $(2)$}{06_endpoint_bit.pdf}{
The proposed one step equation is $4x=3n+1$. Sources $n\equiv1\pmod4$ make $x$ integral.\par
But exact valuation 2 requires the quotient to be odd. The source 5 gives quotient 4, which still contains factors of two. Its actual valuation is 4.\par
The exact sources are $n\equiv1\pmod8$. This final oddness condition explains why a word of total depth $S$ uses modulus $2^{S+1}$ rather than just $2^S$.
}{The final bit distinguishes an exact valuation from a lower bound on valuation.}

\GuideSpread{A cylinder is an arithmetic address}{Refining the initial itinerary of 27}{07_residue_ladder.pdf}{
A word $(a_0,\ldots,a_{k-1})$ is realized exactly when $3^kn+C\equiv2^S\pmod{2^{S+1}}$. Since $3^k$ is odd, it has a unique inverse modulo this power of two. Thus the word defines one odd residue class.\par
The ladder follows four successive prefixes from 27. Each step narrows the arithmetic address. Every row still contains infinitely many positive integers.\par
Having the same coarse residue fixes only the prefix visible at that precision. It does not fix an entire infinite itinerary.
}{Finite precision gives a finite itinerary, with an infinite cylinder of sources.}

\GuideSpread{Composition transforms the earlier offset}{The word $(1)$ followed by the word $(3)$}{08_composition.pdf}{
Let the first block have data $(k,S,C)$ and the second have $(\ell,U,D)$. Their concatenation has offset $3^\ell C+2^SD$.\par
The second block multiplies the earlier offset. It also contributes its own offset at the earlier dyadic scale.\par
For the displayed blocks this gives $3+2=5$, recovering the ordered word $(1,3)$.\par
The resulting affine formula describes an actual orbit only when the source realizes the concatenated word.
}{Transfer blocks compose exactly, including their admissibility conditions.}

\GuideSpread{A local fall can remain above the source}{The first two odd steps from 27}{09_source_floor.pdf}{
The first step rises from 27 to 41. The second falls to 31.\par
The second step is a local decrease, yet it has not fallen below 27. A proof using the original source as its floor cannot reset that floor to 41.\par
The same pattern holds uniformly for $n=16t+11$: $A(n)=24t+17$ and $A^2(n)=18t+13$. The endpoint is below the local peak and above the original source for every integer $t\ge0$.
}{Keep one floor until strict descent below that floor has been proved.}

\GuideSpread{The whole finite orbit from 27}{Forty one odd steps, with the first source descent at 37}{10_orbit27.pdf}{
The curve contains every odd state from 27 through 1. A logarithmic vertical scale keeps both the peaks and the final small values visible.\par
The maximum odd state is 3077 at accelerated step 32. The first state below 27 is 23 at step 37.\par
At that descent $S_{37}=59$ and $C_{37}=1100931843921811423$. The endpoint satisfies $2^{59}\cdot23=3^{37}\cdot27+C_{37}$.\par
Reaching 1 takes 41 odd steps, 70 shortcut steps, or 111 classical steps.
}{A full finite orbit is a certificate for its seed, with exact clock translations.}

\GuideSpread{A contracting coefficient can still miss descent}{A realized 17 step word from 165}{11_contraction_threshold.pdf}{
Here $S=27$, $k=17$, and $2^{27}>3^{17}$. The homogeneous coefficient contracts.\par
Yet $C=1106233681$ exceeds $(2^{27}-3^{17})\cdot165$. The endpoint is 167, above 165.\par
The exact threshold is $\lfloor C/(2^{27}-3^{17})\rfloor=217$. For this word, actual descent requires a realizing source above 217.\par
The source 165 already descends at its first step. This example concerns a finite transfer budget. It is not a least failure candidate.
}{Coefficient crossing is necessary for descent and does not by itself suffice.}

\GuideSpread{Why strict descent is enough}{The strong induction bridge to termination}{12_induction.pdf}{
The example uses $A^4(7)=5$ and $A(5)=1$. Each selected link reaches a smaller positive odd integer.\par
The general bridge is strong induction. Assume every smaller odd source reaches 1. If the current source has any later state below itself, that state reaches 1 by the induction hypothesis. Termination transfers back to the source.\par
Thus Collatz is equivalent to strict descent somewhere along every odd orbit starting above 1. The proof of this implication is already available. The unresolved theorem is the universal descent input.
}{A complete descent theorem needs no additional convergence argument.}

\GuideSpread{No single finite horizon can suffice}{Actual unit valuation prefixes beyond any chosen depth}{13_fixed_horizon.pdf}{
For a proposed horizon $K$, choose $n=2^{K+2}-1$. For every $0\le t\le K$, the actual state is $A^t(n)=2^{K+2-t}3^t-1\ge n$.\par
The curves show exact prefixes for five values of $K$. Each has unit valuations and no early source descent.\par
The formula supplies every horizon, not merely the five drawn cases. A finite table whose rules all prescribe fixed times therefore cannot be a complete descent system.
}{A global certificate may need an unbounded time chosen from its input.}

\GuideSpread{A long run can exit directly to 1}{The exact packet $151,227,341,1$}{14_run151.pdf}{
The source has $j=2$ and $q=19$ in $n=2^{j+1}q-1$. Its first two valuations are 1.\par
The exit relation is $3^3\cdot19=2^9+1$. At the final state 341, $3\cdot341+1=1024=2^{10}$. Thus the final valuation is 10. The word is $(1,1,10)$.\par
The earlier states exceed 151, and the third odd step reaches 1. This certifies both termination and exact first descent time for the seed.
}{Run compression preserves the actual exit and its exact valuation.}

\GuideSpread{Termination can have arbitrarily delayed descent}{Four exact members of a closed seed family}{15_delayed_seeds.pdf}{
Define $q_j=(2^{3^j}+1)/3^{j+1}$ and $n_j=2^{j+1}q_j-1$.\par
The quotient $q_j$ is a positive integer. Factoring $u^3+1$ supplies the induction proving its divisibility.\par
For every $j\ge1$, the first $j$ valuations are 1. The next valuation is $3^j+1$, and that step reaches 1.\par
The table displays $j=1,2,3,4$ using exact decimal integers. These are terminating seeds with exact first descent times $j+1$.
}{Large delay is compatible with termination, so a global proof must allow it.}

\GuideSpread{Lift a seed to an infinite descent progression}{One exact word, unbounded integer sources}{16_cylinder_lift.pdf}{
The word $(1,1,10)$ has total depth 12 and exact cylinder modulus $2^{13}=8192$.\par
For every integer $s\ge0$, $A^3(151+8192s)=1+54s$. The first two states stay above the source. The final target is strictly smaller because the source starts larger and has a much larger progression slope.\par
The plotted values use $s=0$ through 6. A symmetric logarithmic scale keeps the target 1 visible. The uniform formula is proved independently of this plot.\par
The target for $s>0$ is not automatically 1. Its termination requires a further certificate.
}{An infinite descent family is stronger than a list of seeds and weaker than universal termination.}

\GuideSpread{One coarse rule can contain many final valuations}{The progression $4096t+1087$}{17_progression1087.pdf}{
Six exact valuations $1,1,1,1,1,4$ carry the source to $5832t+1549$. Its next numerator is $8(2187t+581)$.\par
Consequently the seventh odd state is the odd part of $2187t+581$. It is at most that positive integer, which is strictly below $4096t+1087$.\par
The red points show that extra divisions sometimes make the exit smaller still. A fixed last valuation is unnecessary for this descent bound.\par
The graph uses $0\le t\le16$. The written inequalities certify all integer parameters $t\ge0$.
}{An upper bound on the odd part can certify a whole progression uniformly.}

\GuideSpread{Read a residue coverage count correctly}{The saved certificate modulo 4096}{18_historical_coverage.pdf}{
The historical atlas records 1804 covered odd classes among 2048, with 244 residual classes. The covered fraction is 88.0859 percent.\par
These squares display that saved aggregate. Their arrangement does not reproduce the residue addresses, and this edition has not rerun the historical coverage computation.\par
Covered classes have selected uniform descent rules. Residual classes are outside that collection. They are not counterexamples.\par
The count is a finite computation result. It is not a Lean theorem enumerating every class and is not a percentage of Collatz proved.
}{A partial certificate has a precise scope even when its coverage is large.}

\GuideSpread{Measure distance from coefficient crossing}{Integer deficit coordinates on a finite prefix}{19_deficit_coordinates.pdf}{
Put $\rho_r=2^{S_r}/3^r$. Permanent coefficient noncrossing means $\rho_r\le1$ at every time. In that domain $d_r=\lfloor r\log_2 3\rfloor-S_r\ge0$.\par
The shaded separation illustrates this difference for the actual first 36 steps from 27. The critical depth is computed by exact integer bit length, not a floating logarithm.\par
At step 37 this orbit crosses and descends. Its finite prefix illustrates the coordinates. It is not evidence that a permanently noncrossing orbit exists.
}{A deficit measures the homogeneous coefficient, not the actual distance above the source.}

\GuideSpread{The rare cylinder bound decays exactly}{Counting words without a random orbit assumption}{20_density_decay.pdf}{
For integers $\ell\ge1$ and $D\ge0$, consider words of length $5\ell$ and depth at most $h=8\ell+D+1$. Sum their cylinder densities:
\[
\varepsilon_\ell=\frac12\sum_{s=5\ell}^{h}\binom{s-1}{5\ell-1}2^{-s}.
\]
Using $z=3/4$ gives
\[
\varepsilon_\ell\le\frac12\left(\frac43\right)^{D+1}\left(\frac{65536}{84375}\right)^\ell.
\]
The exact ratio is below 1. Dots are exact finite rational sums. Lines render the exact rational bounds for $D=0,3$.\par
This counts arithmetic cylinders. It does not assume independent random valuations along an orbit.
}{Rarity follows from exact counting, with a stated word family and depth cap.}

\GuideSpread{Rarity needs a shared height bound}{Why bounded deficit gives a packing contradiction}{21_common_height.pdf}{
This is a proof schematic, not orbit data. Under permanent noncrossing and eventual $d_r\le D$, every late state satisfies $x_r<2^{D+1}(n+r/3)$.\par
The orbit is injective. Thus $R-N$ distinct states must fit below one height $H_R$ proportional to $R$.\par
The rare cylinders have at most $\varepsilon_\ell H_R+W_\ell$ representatives there. Choose $\ell$ so its slope is below the required slope. Letting $R$ grow yields a contradiction.\par
Once deficit is unbounded, this linear height estimate is unavailable. Cylinder rarity by itself cannot finish that branch.
}{The count and the orbit states must be compared below the same height.}

\GuideSpread{Divisibility constrains cycles and mergers}{Necessary equations, with the existence problem visible}{22_divisibility.pdf}{
A nonempty positive cycle satisfies $(2^S-3^k)n=C>0$. The positive gap is necessary. Its existence does not by itself eliminate the cycle equation.\par
For positive sources $n>m$ sharing one word, subtracting their endpoint equations gives $3^k(n-m)=2^S(y-x)$. Coprimality forces $2^S\mid(n-m)$.\par
When $2^S\ge n$, a smaller positive source cannot share that word. This obstructs pulling a smaller endpoint backwards through an unchanged late prefix. An alternative merger must change the word and prove actual realization.
}{An exact necessary relation is useful without being a universal existence theorem.}

\GuideSpread{The two remaining coefficient branches}{What a least possible failure would still need to satisfy}{23_open_branches.pdf}{
If there were a least odd seed failing to reach 1, every later state would stay above that seed.\par
The first branch has some contracting prefix, with $\rho_r>1$. It still must overcome the offset to fall below the original floor.\par
The second branch has permanent noncrossing. Injectivity and the bounded deficit exclusion hold there. The deficit must be unbounded on every tail. This does not say it tends to infinity.\par
Neither branch is closed here. A global argument must force actual descent or contradict every remaining possibility.
}{The unresolved content is a forcing theorem for every odd source above 1.}

\GuideSpread{Three checks to carry away}{Rebuild the exact certificates by hand}{24_reader_checks.pdf}{
For the first line, reconstruct $C_1=1$ and $C_2=5$. Solve the cylinder congruence to obtain $n\equiv11\pmod{16}$.\par
For the second line, verify the valuations $(1,1,10)$. The source and target differ by $8192$ and $54$ from the seed packet.\par
For the third line, start at $5183=4096+1087$. Six prescribed steps give $7381$. Its numerator is $22144=128\cdot173$, so the seventh valuation is 7.\par
Each check has an exact source, word, endpoint, and comparison. The companion monograph supplies the proofs behind their uniform families.
}{Ask which source, which word, which clock, and which floor the certificate preserves.}
